\clearpage
\fancyhead[L]{\small GAUSS AN VON ZACH.}
\hypertarget{native-286-lower}{}
\begin{center}
\textbf{GAUSS AN VON ZACH.}

\hrulefill

Monatliche Correspondenz. Band XV. S. 377. 378. April 1807.

\hrulefill
\end{center}

\begin{flushright}
Braunschweig, 10. März 1807.
\end{flushright}

\par\indent Hiermit habe ich die Ehre, Ihnen meine ersten Beobachtungen der \emph{Pallas} in diesem Jahr zu übersenden. Das schlechte Wetter hat mich bisher daran gehindert, die Beobachtungen fortzusetzen.

\begin{center}\begingroup
\fontsize{9}{11}\selectfont
\setlength{\tabcolsep}{3pt}\renewcommand{\arraystretch}{1.15}
\begin{tabular}{l|l|l|l}
1807 & Mittlere Zeit & \shortstack{Scheinbare\\Rektaszension $\PallasSymbol$} & \shortstack{Scheinbare nördliche\\Deklination $\PallasSymbol$} \\
\hline
Februar 28 & $11^{\mathrm u}\,55^{\mathrm m}\,42^{\mathrm s}$ & $236^\circ\;48'\;6''$ & $8^\circ\;10'\;15''$ \\
März \phantom{0}3 & $12\;\;55\;\;25$ & $237\;\;15\;\;5$ & $8\;\;54\;\;19$ \\
--- \phantom{0}5 & $12\;\;22\;\;2$ & $237\;\;30\;\;41$ & $9\;\;24\;\;52$ \\
\end{tabular}\endgroup\end{center}

\par\indent Die ersten beiden Beobachtungen sind weniger zuverlässig, besonders was die Deklinationen angeht. Die dritte dagegen, die unter günstigeren Umständen gemacht wurde, ist so genau, wie es meine Hilfs%
\SourceLeaf{287}%
mittel erlauben. Zum Vergleich dienten am fünften März zwei Sterne neunter Größe aus der \emph{Hist. cél.}, deren scheinbare Positionen wie folgt angenommen wurden

\begin{center}\begingroup
\fontsize{9}{11}\selectfont
\setlength{\tabcolsep}{3pt}\renewcommand{\arraystretch}{1.15}
\begin{tabular}{l|l}
Rektaszension & Nördliche Deklination \\
\hline
$237^\circ\;19'\;12''$ & $9^\circ\;27'\;50''$ \\
$237\;\;31\;\;13$ & $9\;\;28\;\;8$ \\
\end{tabular}\endgroup\end{center}

\par\indent Nach diesen Beobachtungen gibt meine Ephemeride die Deklinationen genau an. Die Rektaszensionen sind darin um 12 bis 13 Minuten zu groß. Diese Abweichung scheint noch zuzunehmen. Für die Mailänder Beobachtungen der \emph{Ceres} von 1806 danke ich Ihnen sehr. Dadurch habe ich nun auch diese Opposition mit der gewünschten Genauigkeit bestimmt. Allerdings finde ich, dass die Deklination am 4. Januar um 1 Minute zu klein ist und $29^\circ\;21'\;29''3$ betragen sollte (Februarheft 1807, \href{https://archive.org/details/bub_gb_sKAAAAAAMAAJ/page/n164/mode/1up}{S.~159}). Das ist kein Schreibfehler, denn auch die Breite ist zu klein, und Herr \textsc{Carlini} muss sich zufällig auch bei der aus den Elementen berechneten Breite um 1 Minute geirrt haben. Machen Sie diesen geschickten Astronomen doch bei Gelegenheit darauf aufmerksam. Mein letztes Ergebnis stimmt fast genau mit dem von Herrn \textsc{Carlini} berechneten überein. Ich finde die Opposition am 3. Januar 1806 um $12^{\mathrm u}\,20^{\mathrm m}\,12^{\mathrm s}$ mittlerer Zeit in Mailand bei $103^\circ\;1'\;9''4$ wahrer Länge und $6^\circ\;28'\;32''4$ geozentrischer nördlicher Breite. Mit der Verwendung dieses Ergebnisses werde ich nun warten, bis auch die diesjährige Opposition beobachtet worden ist.

